% Source: https://artofproblemsolving.com/community/c3h1715908_junior_regional_olympiad__fbh_2015_grade_7_problem_3
% Stable identity: aops-topic-1715908-q01
Let $AB=AD=a$, $BC=x$, $CD=y$, so $x+y=10$. The diagonal $BD$ is the hypotenuse of both right triangles $BAD$ and $BCD$:

\begin{asy}
usepackage("amsmath");
size(5.5cm);
defaultpen(fontsize(10pt));
real r=sqrt(58)/2;
pair B=(-r,0),D=(r,0),A=(0,r);
real x=10/r; pair C=(x,-sqrt(r^2-x^2));
void ra(pair P,pair U,pair V,real s){draw(P+s*unit(U-P)--P+s*unit(U-P)+s*unit(V-P)--P+s*unit(V-P),linewidth(0.5));}
draw(A--B--C--D--cycle,linewidth(0.9));
draw(B--D,dashed+gray);
ra(A,B,D,0.45);ra(C,B,D,0.45);
// equal-length ticks on AB and AD
pair m1=(A+B)/2,m2=(A+D)/2;
draw(m1+0.22*unit(rotate(90)*(B-A))--m1-0.22*unit(rotate(90)*(B-A)));
draw(m2+0.22*unit(rotate(90)*(D-A))--m2-0.22*unit(rotate(90)*(D-A)));
label("$A$",A,N);label("$B$",B,W);label("$C$",C,S);label("$D$",D,E);
label("$a$",m1+0.6*unit(rotate(-90)*(B-A)));label("$a$",m2+0.6*unit(rotate(90)*(D-A)));
label("$x$",(B+C)/2,SW);label("$y$",(C+D)/2,SE);
label("\mbox{}",(min(currentpicture,user=true)+max(currentpicture,user=true))/2,black); // keeps the text after the figure black when LaTeX typesets the labels inline
\end{asy}

\[
BD^2=a^2+a^2=2a^2\qquad\text{and}\qquad BD^2=x^2+y^2 .
\]
The area is the sum of the areas of the two right triangles:
\[
[ABCD]=\dfrac{a^2}{2}+\dfrac{xy}{2}=\dfrac{2a^2+2xy}{4}=\dfrac{x^2+y^2+2xy}{4}=\dfrac{(x+y)^2}{4}=\dfrac{100}{4}=\boxed{25}.
\]
(The answer does not depend on how the sum $10$ is split between $x$ and $y$; the figure shows $x=7$, $y=3$.)
